\documentclass[11pt]{article}
\usepackage[margin=1in]{geometry}
\usepackage{amsmath,amssymb,booktabs,longtable,hyperref}
\newcommand{\pa}{pseudo-Anosov}
\title{A homological lower bound for twenty-one genus-three strata}
\author{Computational proof record}
\date{10 October 2026}
\begin{document}\maketitle
\begin{abstract}
For each of the twenty-one strata listed below, every orientation-preserving
\pa{} homeomorphism of a closed orientable genus-three surface has dilatation
at least $c$, the largest real root of
$C(X)=X^6-X^4-X^3-X^2+1$. The statement includes every connected component and
every permutation of zeros of the same order. It concerns actual maps, not
the global minimum of an abstract train-track automaton. It does not identify
any stratum's exact minimum and does not prove the unrestricted genus-three
minimum. The topological argument is written here; the attached programs
check finite exact arithmetic and provenance, not a formalization of topology.
\end{abstract}

\section{Scope and the necessary homological spectra}
The Abelian strata are $H(4),H(3,1),H(2,2),H(2,1,1),H(1,1,1,1)$.
The quadratic strata in this theorem have orders
\begin{align*}
 b=0:&\quad (8),(6,2),(4,4),(4,2,2),(2,2,2,2);\\
 b=2:&\quad (7,1),(6,1,1),(5,3),(5,2,1),(4,3,1),\\
 &\quad(4,2,1,1),(3,3,2),(3,2,2,1),(2,2,2,1,1);\\
 b=4\text{, all orders odd}:&\quad(3,3,1,1),(5,1,1,1).
\end{align*}
Here $b$ is the number of odd quadratic orders; $Q$ means the foliation is
nonorientable. A square quadratic differential belongs to the Abelian case.
No assertion about $Q(4,1,1,1,1)$, $Q(3,2,1,1,1)$, $Q(2,2,1,1,1,1)$,
$Q(3,1,1,1,1,1)$, $Q(2,1,1,1,1,1,1)$, or $Q(1^8)$ is made here.

For an orientable foliation choose its global orientation. The sign of an
orientation-preserving map on both foliations is a single
$\varepsilon\in\{+1,-1\}$. Its genus-three homology action has simple
strictly dominant eigenvalue $\varepsilon\lambda$ by the orientable-foliation
homology theorem \cite[Theorem 2.2]{LT}. Multiplying the action by
$\varepsilon$ gives an integral symplectic matrix and a reciprocal monic
degree-six characteristic polynomial with constant coefficient one and
simple strictly dominant eigenvalue $\lambda>1$.

For a nonorientable quadratic foliation, take its connected orientating
double cover \cite[Construction 2.1]{LT}, with deck involution $\tau$.
Riemann--Hurwitz gives cover genus $5+b/2$. The invariant homology has rank
six and carries the base action; the anti-invariant homology has rank $4+b$.
The kernel of $1+\tau_*$ in integral homology is an invariant lattice with
nondegenerate symplectic form over $\mathbb Q$. Thus its characteristic
polynomial is integral, monic and reciprocal, with constant coefficient one.
Choose the foliation-preserving lift $F$ of the map. The simple strictly
dominant homological eigenvalue $\lambda$ lies in this anti-invariant space,
because the oriented foliation changes sign under $\tau$. All base
eigenvalues have modulus strictly less than $\lambda$.

Set $u=701/500$; exact Sturm isolation gives $c<u$.
Suppose, for contradiction, $1<\lambda<c$. All the foregoing necessary
characteristic polynomials then have every root strictly inside $|X|<u$.
No train-track coverage assumption is used in this argument.

\section{The complete small-house catalogs}
For rank four write $P(X)=X^2 Q(X+X^{-1})$ with
$Q(Y)=Y^2+aY+b$. Each root $y$ of $Q$ is $z+z^{-1}$ for a root $z$ of
$P$. Reciprocity gives $u^{-1}<|z|<u$, whence
$|y|<T=u+u^{-1}$. Consequently $|a|<2T<5$ and $|b|<T^2<5$.
The containing box $-4\leq a,b\leq4$ consists of exactly 81 polynomials
$[1,a,b+2,a,1]$. The full box is tested by three independent exact unit-disk
criteria after rational scaling: Schur recursion, a Schur--Cohn Toeplitz
test, and a Cayley transform with a Hurwitz test. Every stable polynomial
has zero real roots in $(1,u)$ by exact Sturm counts (factors $X-1$ are
removed before counting). Thus all five unbranched quadratic strata are
excluded under $\lambda<c$.

For degree six the complete catalog is the previously independently
reconstructed 75-polynomial small-house packet. It is derived by a
containing coefficient box for $P(X)=X^3Q(X+X^{-1})$, not a positivity
restriction on traces. The original cubic reconstruction and the direct
15,093-element coefficient box are independently compared in job 9407;
all three disk tests agree. Exactly three catalog polynomials have a real
root in $(1,u)$:
\[
 C(X),\qquad W(X)=X^6-3X^3+1,\qquad
 A(X)=X^6+X^5-X^4-3X^3-X^2+X+1.
\]
$C$ has no real root in $(1,c)$. $W$ is a polynomial in $X^3$ and its
positive largest root has distinct conjugates of the same modulus, violating
strict dominance. Hence any putative actual map with $\lambda<c$ in the
degree-six cases has characteristic polynomial $A$. Its traces begin
$(-1,3,5,-1,9,3)$; all 32 traces used below are checked by both Newton
identities and exact companion-matrix powers.

For the two all-odd $b=4$ strata, positivity of regular point counts implies
$t_-(k)\geq0$. We use the weaker containing bounds
$-2\leq t_-(k)\leq\lfloor 8u^k\rfloor$ for $1\leq k\leq4$.
Newton divisibility reconstructs all 8,719 distinct reciprocal polynomials.
The independent rectangular trace reconstruction agrees with the recursive
reconstruction; all three disk tests leave exactly 120 stable polynomials,
25 with a real root in $(1,u)$, and exactly 20 whose largest positive root
is strictly below $c$. The other five cannot represent $\lambda<c$.
Each of the 20 is either excluded by the exact strict-dominance factor
test or subjected to every necessary zero-action and base-homology model
below. Passing a necessary condition never asserts realizability.

\section{Abelian zero actions, including a reversed foliation}
An Abelian zero of order $m$ has $p=2m+2$ prongs. Zeros of different orders
cannot be permuted. For a cycle of $\ell$ zeros of a given order, the total
return rotation $r\pmod p$ has parity zero when $\varepsilon=+1$, and
parity $\ell\pmod2$ when $\varepsilon=-1$. This follows because consecutive
prongs alternate their signed foliation orientation. Every allowed total
rotation is retained. At iterate $k$ divisible by $\ell$, the cycle's
singular Lefschetz contribution is
\[
 \ell\begin{cases}1-p,&r(k/\ell)=0\pmod p,\\1,&\text{otherwise}.\end{cases}
\]
These are the standard local-index formulas \cite[Proposition 2.6 and
Corollary 2.7]{LT}; nonfixed zeros contribute zero. Let their total be
$L_{\rm sing}(k)$. A regular fixed point has index $-\varepsilon^k$.
Since $L(f^k)=2-\varepsilon^k t_A(k)$, its regular fixed-point count is
\[
 R(k)=t_A(k)+\varepsilon^k\bigl(L_{\rm sing}(k)-2\bigr).
\]
For each of the five Abelian strata, every cycle type and return rotation
of each sign is enumerated. Independently the program runs every labelled
zero permutation and every allowed per-vertex rotation, takes literal
powers on the prongs, and compares the full 32-index vectors with the
cycle model. The sets agree exactly. Every $A$-model fails either
$R(k)\geq0$ or the nonnegative integral ordinary orbit condition
\[
 O(n)=\frac{R(n)-\sum_{d\mid n,\,d<n}dO(d)}{n}\in\mathbb Z_{\geq0}
 \quad(1\leq n\leq32).
\]
As positive controls, compatible $C$ zero-action models are retained, and
the two genuine toral automorphisms $\pm\begin{pmatrix}2&1\\1&1\end{pmatrix}$
give $R(k)=|\det(f_*^k-I)|$ independently. A corrupted period-two count
is rejected. Compatible singular data alone do not identify an attaining
map. Thus no Abelian map has dilatation below $c$.

\section{Quadratic zero actions and base coupling}
A quadratic zero of order $a$ has $p=a+2$ base prongs. We retain all cycles
of zeros of the same order and all return rotations. At each iterate keep
the \emph{correlated} vector $(S(k),Z(k),E(k))$: $S$ is the singular
contribution in $L(F^k)-L(\tau F^k)=-2t_-(k)$, $Z$ the number of fixed
base zeros, and $E$ the sum of $\pm1$ return sheet signs at fixed even
zeros. Even a shared $S$-vector is not enough to discard distinct $Z,E$.
For an odd zero, its branched lift has $2p$ prongs; the positive-foliation
lift of a base rotation $r$ is rotation $(p+1)r\pmod{2p}$.
For an even zero use its two lifts and the parity of the rotation.
Literal powers of these finite prong and sheet actions compute both local
indices. Their difference gives $S$, while their average gives
\[
 \operatorname{Ind}_{\rm base}=Z-S-E.
\]
For a branched point, the average identity also follows from the coordinate
$z=w^2$: $f(w^2)-w^2=(F(w)-w)(F(w)+w)$, so twice the base winding number
is the sum of the two lift winding numbers. Here $\tau F^k$ is meant;
it is not generally $(\tau F)^k$.

Regular points with preserved foliation orientation contribute $-1$ on
the base and two $-1$ indices to $F^k$; those with reversed orientation
contribute $+1$ and instead lift to $\tau F^k$. Thus
\[
 R(k)=t_-(k)-S(k),\qquad
 R_-(k)=\frac{2-t_+(k)-Z(k)+t_-(k)+E(k)}2.
\]
In addition to ordinary nonnegative integral orbit counts, negative-sign
regular orbits satisfy
\[
 O_-(n)=\frac{R_-(n)-
 \sum_{d\mid n,\,d<n,\,n/d\ {
m odd}}dO_-(d)}n
 \in\mathbb Z,\qquad 0\leq O_-(n)\leq O(n).
\]
Only odd returns preserve a negative sign; using every divisor here would
be incorrect. Genuine torus quotient controls, direct residue-class point
counts, and a negative-sign two-cycle check in the accepted generic module
test these identities and distinguish that error.

All cycle/rotation vector sets in this proof are independently compared
with all labelled permutations and all per-vertex rotations and literal
prong/sheet powers. Products across distinct zero orders are full products,
not a minimum-dominance selection. For each of the nine $b=2$ strata,
every $A$-model already violates the ordinary necessary orbit conditions.
This excludes $\lambda<c$ for all nine.

For the two all-odd $b=4$ strata the ordinary conditions alone are
insufficient. Each strict-dominance-admissible polynomial among the 20 is
therefore coupled to \emph{all} 75 base polynomials. For its isolated
dominant root $\lambda\in(l,h)$, a base polynomial is retained only if all
its roots have modulus $<h$; this is a conservative necessary test, not a
claim that every retained base spectrum has modulus $<\lambda$.
All three exact disk tests agree. Exact greatest common divisors reject a
base sharing $\lambda$ or $-\lambda$. Every remaining base polynomial and
every complete correlated zero-action model is tested through period 32.
All models fail the ordinary or coupled necessary conditions. This excludes
both $Q(3,3,1,1)$ and $Q(5,1,1,1)$, including all their components.

\section{Independent arithmetic acceptance and its limits}
The diagnostic is \texttt{genus3\_bulk\_homology\_lower\_diagnostic\_v2\_20261010.py}.
Its original version correctly stopped when it incorrectly expected the
ordinary conditions alone to exclude every all-odd rank-eight model.
Its original source, partial packets, log and terminal metadata are kept.
The new version retains the surviving models and adds the full base
coupling; exact source reversal restores the original. No proof credit is
taken from the failed version.

The final acceptance program freshly replays the complete degree-six
independent arithmetic, the rank-eight recursive catalog, the independently
reconstructed rank-eight catalog, and the complete new diagnostic. Every
full report equals its original apart from elapsed time, and all detail
packets agree byte for byte. It freezes the manuscript, primary paper,
algorithm sources and dependencies, and complete arithmetic packets with
SHA-256 manifests and read-only snapshots. The acceptance explicitly binds
the twenty-one universal real-map lower bounds to this written argument.
It grants no new automaton base coverage, automaton global minimum,
stratum exact minimum, or unrestricted global genus-three minimum.

\begin{thebibliography}{9}
\bibitem{LT} E. Lanneau and J.-L. Thiffeault,
\emph{On the minimum dilatation of pseudo-Anosov homeomorphisms on surfaces
of small genus}, arXiv:0905.1302, Construction 2.1, Theorem 2.2,
Proposition 2.6, Corollary 2.7.
\end{thebibliography}
\end{document}
